\documentclass[10pt]{article}

\usepackage[utf8]{inputenc}

\usepackage[T1]{fontenc}

\usepackage{amsmath}

\usepackage{amsfonts}

\usepackage{amssymb}

\usepackage[version=4]{mhchem}

\usepackage{stmaryrd}

\usepackage{graphicx}

\usepackage[export]{adjustbox}

\graphicspath{ {./images/} }

\usepackage{CJKutf8}

\usepackage{bbold}



\begin{document}

ски - векторным полем), то говорят о Ф. т. автономной системы (поля), используя это название и тогда, когда решения системы определены не для всех значений $t$. Прилагательное «фазовая» нередко опускают.



Когда состояние не зависит от $t$, Ф. т. сводится к точкө - равновесия положению, а при периодич. зависимости от $t$ получается замкнутая Ф. т. (поэтому часто говорят о «периодич. Ф. т.»), что включает и предыдущий случай (однако, говоря о замкнутой Ф. т., часто подразумевают, что она не сводится к точке). Незамкнутые Ф. т. в общем случае могут быть весьма разнообразны; они классифицируются с различных точек зрения в топологической динамике. Точка $w$ незамкнутой Ф. т. разбивает еө на две части - п о л ожи т е льн ую и о т р ица те Јь н ую п о л ут р а е к т о р и и. Они представляют состояния, соответствующие $t \geqslant 0$ и $t \leqslant 0$, если при $t=0$ система имеет состояние $w$. (Последнее определение формально применимо и к замкнутой Ф. т., но для нее обе полутраектории совпадают с ней самой.)



Иногда под $\Phi$. т. понимают не просто линию (как множество точек) или ориентированную линию (на к-рой выделено направление, отвечающее изменению состояний с ростом $t$ ), но линию, параметризованную в процессе происходящего в системе движения фазовой точки по этой линии. Такое словоупотребление отчасти связано с отсутствием удобного общепринятого названия для этой параметризованной линии. Правда, если динамич, система описывается системой дифференциальных уравнений, речь идет просто о решении последней; но такое название не годится в общем случае, когда динамич. система трактуется как группа преобразований $\left\{S_{t}\right\}$ фазового щространства. (Функцию $t \rightarrow S_{t} w$ с фиксированным $w$ иногда наз. "движением», но в математике «движения» обычно являются шреобразованиями всего пространства.) Д. В. Аносов.



ФАЗОВОГО РАВНОВЕСИЯ ДИАГРАММА - проекция на плоскость каких-либо двух термодинамич. переменных тех областей поверхности равновесных состояний в пространстве полного набора термодинамиq. переменных, к-рые соответствуют $n$-фазным, $n \geqslant 2$, состояниям термодинамич. системы. В случае однокомпонентной системы участки этой поверхности представляют собой цилиндрич. поверхности и гроектируются на плоскость $(p, \theta)$ (давлөние-температура) в виде линии, общий вид уравнения к-рой - равенство химич. потенциалов различных ф্রаз $\mu_{1}(p, \theta)=\mu_{2}(p, \theta)-$ может быть в случае фазового перехода 1-го рода записан в форме дифференциального уравнения Клапейрона - Клаузиуса



\[

d p / d \theta=q_{12} / \theta\left(v_{2}-v_{1}\right),

\]



где $q_{12}$ - скрытая теплота перехода, $v_{1}$ и $v_{2}$ - удельные объемы для первой и второй фаз. Трехфазное состояние изображается точкой, наз. т р о й н о й т 0 ч к о й.



\begin{center}

\includegraphics[max width=\textwidth]{2023_07_09_ea0bd083229ceac547b6g-1}

\end{center}



ФАЗОВОЕ ПРОСТРАНСТВО - совокупность всевозможных мгновенных состояниӥ физич. (в широком смысле слова) системы, снабженная апределенной структурой в зависимости от изучаемой системы и рассматриваемых вопросов. Ф. п. наз. также более конкретный объект - пространство (множество с надлежащей структурой), элементы к-рого (ф а з о в ы е т о ч к и) представляют (условно изображают) состоявия системы (напр., фазавая плоскость). С математич. точки зрения эти объекты изоморфны, поэтому часто не делают различия между состояниями и изображающими их фазовыми точками.



Математич. формализация понятия «системы» того или иного типа обычно включает в качестве существенной части определение соответствующего Ф. п. (или класса Ф. п.), что отражает важность понятия состояния системы. Эволюция системы (т. е. изменение со временем ее состояний) может быть строго детерминированной (тогда она описывается группой или полугруппой $\left\{S_{t}\right\}$ греобразованиї Ф. п.: состояние $w$ за время $t$ переходит в $S_{t} w$ ) или иметь вероятностный характер (случайные прочессы). В первом случае тожке бывает нужно рассматривать статистич. состояния системы; для классич. (неквантовых) систем они описываются посредством вероятностных распределений на Ф. П. Указания, определяющие эволюцию, составляют другую существенную часть формального определения «системы».



В классич. случае диффоренцируемой динамической системы (что включает основные системы, рассматриваемые в аналитич. мөханике и классич. статистич. физике) Ф. п.- дифберенцируемое многообразие $M$ (быть может, с особенностями и (или) краем). Если динамич. система задается автономной системой обыкновенных дифференциальных уравнений, то говорят o Ф. п.' последней. В әтом случае $M$ является той областыо в евклидовом пространстве, где определены правые части автономной системы. В подобной ситуации термин «Ф. п.» употребляют и тогда, когда решения не определены при всех $t$. Дополнительно на $M$ может быть задана инвариантная мера (классически - задаваемая плотностью) или симплектическая структура (условие сох ранения последней под действием потока характеризует гамильтоновь системь). В частности, в динамике для системы с голономными идеальными связями, не зависящими явно от времени, Ф. п. является касательным или кокасательным расслоением нек-рого многообразия - к о н ф্\begin{CJK}{UTF8}{mj}丶\end{CJK} и г у р а ц и о н н ог о п р о с т р а н т в а. Точка последнего задае' (рас)положение (конфигурацию) системы, касательный вектор описывает скорость движения системы (скорость изменения конфиигурации), а кокасательный - импульс системы.



В других разделах теории динамич. систем Ф. п. имеет структуру тоголотич. пространства (в топологической динамике), измеримого пространства или (чаще) пространства с мерой (в эргодической теории). В квантовой механикө Ф. п.- комплексное гильбертово пространство (впрочем, для квантовой системы, имеющей классический аналог, под Ф. п. часто подразумевают Ф. П. этого аналога). В теории случайных процессов Ф. п.- то измеримое пространство (часто с дополнительной топологической, диффференцируемой или векторной структурой), в к-ром принимает значения процесс. Здесь специально говорят о Ф. п. тогда, когда оно в каком-то смысле нетривиально. Так нередко бывает в теории марковских процессов, тогда как для часто встречающихся процессов с числовыми значениями Ф. п. сводится просто к $\mathbb{R}$ со стандартной структурой, и споциально говорить о нем не приходится. (Надо иметь в виду, что если стационарный в узком смысле случайный процесс интерпретируется как динамич. система, то Ф. п. последней не совпадает $\begin{array}{ll}\text { с Ф. п. прогесса.) } & \text { Д. В. Аносов. }\end{array}$



ФАЗОВОЙ СКОРОСТИ ВЕКТОР - вектор $f(x)$, исходящий из точки $x$ фазового пространства $G$ автономной системы



\[

\dot{x}=f(x), t \in C^{1}(G), G \subset \mathbb{R}^{n} .

\]



Пусть $\Gamma$ - фазовая траектория этой системы, проходящая через точку $\xi \in G$; если $f(\xi) \neq 0$, то Ф. с. в. $f(\xi)$ касается кривой Г в точке $\xi$ и представляет собой мгновенную скорость движения изображающей точки системы по траектории $\Gamma$ в момент грохождения положкеня $\xi \in \Gamma$. Если $f(\xi)=0$, то точка $\xi \in G$ является равновесия положением.



Jlum.: [1] П о н т р я г и н.Л. С., Обыкновенные дифференциальные уравнения, 5 изд., М., 1983.



H. X. Розов.





\end{document}